\theory{Index theory}{How can the difference between solutions and obstructions of a differential equation be computed from topology?}{The index of a Fredholm operator is the dimension of its kernel minus the dimension of its cokernel. The Atiyah--Singer theorem identifies this analytic integer with a topological expression.}{Elliptic partial differential equations, geometry, topology, heat flow, quantum physics, and moduli spaces.}{The Fredholm index turns an analytic kernel-minus-cokernel difference into a stable topological quantity under perturbation.}

\paragraph{The problem and history}
An equation $D u=f$ may have solutions, many solutions, or no solution for some right-hand sides. The kernel contains solutions of $Du=0$. The cokernel records the independent compatibility conditions that $f$ must satisfy. For a Fredholm operator, both spaces are finite-dimensional and
\[
\operatorname{ind}(D)=\dim\ker D-\dim\operatorname{coker}D.
\]
Israel Gel'fand asked whether this integer could be expressed using topological data. Michael Atiyah and Isadore Singer answered yes in the 1960s \cite{index_theory_ref1,index_theory_ref2}.

\end{historylist}
\section*{3. Theory}
\subsection*{Core theory}
\paragraph{Why the difference is stable}
The dimensions of the kernel and cokernel can jump when an operator is perturbed. Their difference often does not. On a circle, the operator $D=d/dx-\lambda$ on periodic functions has a nonzero kernel only for special values of $\lambda$, but its index remains zero because the adjoint has a matching kernel.

Elliptic operators on compact manifolds are Fredholm. Ellipticity means that the highest-order part of the operator has no hidden direction in which information can propagate without control. Compactness prevents infinitely many independent low-energy solutions.

\paragraph{Symbols and the Atiyah--Singer formula}
The leading derivative part of an elliptic operator is its symbol. The symbol is a bundle map on the cotangent bundle and is invertible away from the zero covector. Its homotopy class defines a K-theory element. The analytic index depends only on this topological symbol class.

The Atiyah--Singer index theorem states
\[
\operatorname{ind}_{\mathrm{an}}(D)=\operatorname{ind}_{\mathrm{top}}(D).
\]
The topological side is built from characteristic classes of the manifold and bundles. One can therefore count analytic zero modes without solving the differential equation explicitly.

\paragraph{Major special cases}
The Riemann--Roch theorem computes the dimension difference between holomorphic sections and obstructions. The Hirzebruch--Riemann--Roch theorem extends it to complex manifolds. The signature theorem expresses the signature of the intersection form using the $L$-class. The Chern--Gauss--Bonnet theorem computes the Euler characteristic from curvature. The Dirac operator yields the Â-genus and spin-geometric results \cite{index_theory_ref3}.

\paragraph{Proof viewpoints and applications}
Atiyah and Singer first used K-theory and a push-forward construction, reducing the problem to spheres. Cobordism gives another reduction to simple generators. A heat-equation proof writes the index as a difference of traces of heat operators; the small-time asymptotic expansion reveals the characteristic classes. Supersymmetry explains cancellations in the heat expansion.

Index theory counts zero modes in gauge theory, instantons, and quantum field theory. It gives dimensions of moduli spaces, integrality constraints such as Rokhlin-type results, and topological information from elliptic boundary problems. The Atiyah--Patodi--Singer theorem adds boundary correction terms involving eta invariants.

\paragraph{What the theory depends on and where it is useful}
Index theory depends on functional analysis, elliptic PDE, differential geometry, topology, K-theory, and characteristic classes. It helps Hodge theory identify harmonic forms, geometry compute invariants, and physics count stable states.

The index does not usually give the separate dimensions of the kernel and cokernel. On noncompact spaces or for non-elliptic operators, finite-dimensionality and stability may fail; boundary conditions must then be chosen carefully.

\section*{4. Important theorems and results}
\begin{enumerate}[leftmargin=*,itemsep=0.35em]
\item \textbf{Atiyah--Singer index theorem.} The analytic and topological indices of an elliptic operator on a compact manifold are equal.

\item \textbf{Riemann--Roch theorem.} A holomorphic dimension difference is computed by degree and genus.

\item \textbf{Chern--Gauss--Bonnet theorem.} The Euler characteristic is the integral of a curvature characteristic form.

\item \textbf{Hirzebruch signature theorem.} The signature of a compact oriented manifold is an integral of its $L$-class.

\item \textbf{Atiyah--Patodi--Singer theorem.} Elliptic index formulas on manifolds with boundary include a boundary spectral correction \cite{index_theory_ref4}.

The index is an integer that measures an analytic imbalance, such as the difference between the dimensions of a kernel and a cokernel. Atiyah--Singer says that this analytic count equals a topological expression, so it can be computed without solving the differential equation directly. Riemann--Roch, Gauss--Bonnet, and the signature theorem are important special cases; Atiyah--Patodi--Singer shows what changes when a boundary contributes spectral information.
\end{enumerate}

\theoryapplications
\theoryconnections{index theory}{%
\item Functional analysis supplies operators and spectra, differential geometry supplies elliptic differential operators, and topology supplies characteristic classes and K-theory.
\item Measure and integration theory provide the analytic index, while homological algebra organizes the associated complexes.
}{%
\item Index theory helps compute analytic quantities from topological data in geometry and mathematical physics.
\item It explains why a problem involving infinitely many analytic modes can have a finite, stable integer answer.
}
\section*{Glossary}
\begin{description}[leftmargin=*,style=nextline,itemsep=0.3em]
\item[\textbf{Fredholm operator.}] A bounded linear operator between Banach spaces whose kernel and cokernel are both finite-dimensional, so that the difference of their dimensions (the index) is well-defined.
\item[\textbf{Kernel.}] The subspace of inputs mapped to zero by a linear operator; in index theory it represents independent solutions of $Du=0$.
\item[\textbf{Cokernel.}] The quotient of the target space by the image of a linear operator; it measures independent obstruction conditions that a right-hand side must satisfy.
\item[\textbf{Index.}] The integer $\dim\ker D - \dim\operatorname{coker} D$ of a Fredholm operator; it is stable under small perturbations and, by the Atiyah--Singer theorem, computable from topology.
\item[\textbf{Elliptic operator.}] A differential operator whose highest-order part is invertible away from the zero covector, ensuring good analytic behavior on compact manifolds.
\item[\textbf{Symbol.}] The leading-derivative part of an elliptic operator, viewed as a bundle map on the cotangent bundle; its homotopy class determines the index.
\item[\textbf{Characteristic class.}] A cohomology class associated to a vector bundle that encodes topological information; the Atiyah--Singer formula expresses the index as an integral of characteristic classes.
\item[\textbf{K-theory.}] A generalized cohomology theory built from stable equivalence classes of vector bundles.
\item[\textbf{Riemann--Roch theorem.}] Computes the index of the $\bar\partial$-operator on a complex manifold as a difference of holomorphic dimensions.
\item[\textbf{Chern--Gauss--Bonnet theorem.}] Expresses the Euler characteristic as the integral of a curvature characteristic form.
\item[\textbf{Eta invariant.}] A spectral invariant of an elliptic operator on a manifold with boundary.
\end{description}
\section*{7. Sample Questions}
\emph{Exercise reference:} \cite{index_theory_ref1}. The cited edition is the source for the exercise set; consult its Exercises section for the official statement and numbering.
\begin{enumerate}[leftmargin=*,itemsep=0.9em]
\item \textbf{Exercise 1.} Compute the Euler characteristic of $S^2$ using the Chern--Gauss--Bonnet theorem.

\emph{Solution.} The Chern--Gauss--Bonnet theorem states $\chi(M)=\frac{1}{2\pi}\int_M K\,dA$ for a closed oriented surface, where $K$ is the Gaussian curvature. For $S^2$ with the round metric of radius $1$, $K=1$ everywhere and $\operatorname{Area}(S^2)=4\pi$. Therefore $\chi(S^2)=\frac{1}{2\pi}\cdot 4\pi=2$.

\item \textbf{Exercise 2.} Compute the index of the operator $D=\frac{d}{dx}:\Omega^0(S^1)\to\Omega^1(S^1)$ on the circle (acting on smooth functions to produce exact 1-forms).

\emph{Solution.} The kernel of $D$ consists of functions with $df=0$, i.e., constant functions: $\dim\ker D=1$. The cokernel is $H^1_{\mathrm{dR}}(S^1)\cong\mathbb{R}$, so $\dim\operatorname{coker}D=1$. The index is $\operatorname{ind}(D)=1-1=0$.

\item \textbf{Exercise 3.} State the Riemann--Roch theorem for a compact Riemann surface of genus $g$, applied to the trivial line bundle $\mathcal{O}$.

\emph{Solution.} Riemann--Roch states $\ell(D)-\ell(K-D)=\deg(D)-g+1$, where $D$ is a divisor, $K$ a canonical divisor, $\ell(D)=\dim H^0(\mathcal{O}(D))$, and $\ell(K-D)=\dim H^1(\mathcal{O}(D))$. For $D=0$: $\ell(0)-\ell(K)=0-g+1=1-g$. Since $\ell(0)=1$ (only constants) and $\ell(K)=g$ (the space of holomorphic $1$-forms), we get $1-g=1-g$, a consistency check.

\item \textbf{Exercise 4.} Verify the Atiyah--Singer formula for the $\bar\partial$-operator on a compact Riemann surface of genus $g$.

\emph{Solution.} The index of $\bar\partial:L^2(\Omega^{0,0})\to L^2(\Omega^{0,1})$ is $\dim\ker\bar\partial-\dim\operatorname{coker}\bar\partial=g-1$. The analytic index counts holomorphic functions minus obstructions. The topological side, by Riemann--Roch, gives the same: $\frac{1}{2}(\deg(T^{1,0})-\deg(K))=\ldots$ The $\hat{A}$-genus on a surface of genus $g$ equals $1-g$ (using the standard formula). Both sides agree, confirming the theorem.

\item \textbf{Exercise 5.} Compute the signature of $S^2\times S^2$ using the intersection form.

\emph{Solution.} $H_2(S^2\times S^2;\mathbb{Z})\cong\mathbb{Z}^2$, generated by $a=[S^2\times\{pt\}]$ and $b=[\{pt\}\times S^2]$. The intersection form has matrix $\begin{pmatrix}a\cdot a&a\cdot b\\b\cdot a&b\cdot b\end{pmatrix}=\begin{pmatrix}0&1\\1&0\end{pmatrix}$. The eigenvalues are $+1$ and $-1$, giving one positive and one negative direction. The signature is $\sigma=1-1=0$.

\end{enumerate}
\section*{8. References}
\addcontentsline{toc}{section}{References}

\begin{thebibliography}{9}
\bibitem{index_theory_ref1} Michael F. Atiyah and Isadore M. Singer, \emph{The Index of Elliptic Operators I}, Annals of Mathematics, 1968.
\bibitem{index_theory_ref2} Michael F. Atiyah, \emph{K-Theory}, 2nd ed., Westview Press, 1989.
\bibitem{index_theory_ref3} H. Blaine Lawson and Marie-Louise Michelsohn, \emph{Spin Geometry}, Princeton University Press, 1989.
\bibitem{index_theory_ref4} Richard B. Melrose, \emph{The Atiyah--Patodi--Singer Index Theorem}, A K Peters, 1993.
\end{thebibliography}
